\documentclass[11pt,a4paper]{article}
\usepackage{amsmath}
\usepackage[margin=1in]{geometry}
\title{A Simple Performance Evaluation of a Classifier}
\author{Belan Bapir\\Department of Computer Engineering\\Tishk International University}
\date{December 9, 2026}
\begin{document}
\maketitle
\begin{abstract}
This teaching example explains how to calculate classification accuracy using hypothetical data. It also demonstrates the structure of a scientific paper and automatic cross-references.
\end{abstract}
\section{Introduction}
\label{sec:intro}
Classification assigns samples to categories. Section~\ref{sec:method} explains the evaluation method, and Section~\ref{sec:results} presents an illustrative result.
\section{Methodology}
\label{sec:method}
\subsection{Evaluation Data}
Assume a classifier is tested on 100 samples, of which 85 are classified correctly.
\subsection{Performance Metric}
Accuracy is defined as
\begin{equation}
\label{eq:accuracy}
A=\frac{N_{\mathrm{correct}}}{N_{\mathrm{total}}}\times 100\% .
\end{equation}
\section{Results}
\label{sec:results}
Using Equation~\eqref{eq:accuracy}, the illustrative accuracy is
\begin{equation}
\label{eq:result}
A=\frac{85}{100}\times 100\%=85\%.
\end{equation}
\section{Conclusion}
Equation~\eqref{eq:result} shows that the classifier correctly identifies 85 out of 100 samples.
\section*{Acknowledgements}
Prepared as a teaching example for the TIU national LaTeX workshop. The numerical values are hypothetical, not research findings.
\end{document}
